\chapter{Getting started}\label{ch:start}

\section{The data}\label{sec:data}

A fit needs three things in Stata: a network in memory, a dataset with one row per actor that holds the outcome (and the covariates), and the command. The outcome is a numeric variable with the values 0 and 1; missing values are allowed (Chapter~\ref{ch:model}). The network is the one named in the command, or the current network.

How the rows of the dataset are matched to the actors of the network depends on the dataset. When it carries the node labels (the variable \texttt{\_nwnode}, as after \texttt{nwload} with the option \texttt{xvars}, or after \texttt{nwwebuse}), the rows are matched to the actors by label, in any row order, and a mismatch (an actor with no row or two rows) is an error. Otherwise the dataset must have one row per actor, in the order of the network. For a two-mode network, see Section~\ref{sec:twomode}.

Covariates are variables of the same dataset. Numeric covariates are used as they are, so \emph{centre them} (a distance matrix is also not centred); categorical effects such as \opt{same()} compare the values of the variable, which may be numeric or string.

\section{A first model}\label{sec:first}

The karate club of Zachary~\cite{zachary1977} has 34 members and 78 friendship ties; the outcome is the faction the member took after the club split (\texttt{faction}, 18 of 34 with the value 1). The edges and attributes are those that come with \alaamee{}. We load the network with \texttt{nwfromedge}, then the attributes:
\showlog{karate_data}
The first model has the three effects of the classical example: the number of actors in the faction (\opt{density}, the intercept), the sum of their degrees (\opt{activity}) and the number of ties between two members of the faction (\opt{contagion}):
\showlog{karate_sa}
The header says that stochastic approximation was used, that the outcome has 18 ones among 34 actors, and that the procedure was run twice until its convergence statistic was within $0.1$ in absolute value. The table is the usual coefficient table of a Stata estimation command; the standard errors come from the covariance of the simulated statistics (Section~\ref{sec:sa}). Members of the faction tend to be tied to each other (contagion 1.48, $z=3.65$), after allowing for the fact that members with many ties are less likely to belong to it (activity $-0.84$). The intercept is not significant, which is common and not a concern: it only sets the base rate.

Because the algorithm simulates, the numbers depend on the random-number seed. \opt{seed(\#)} sets it; without it the seed is drawn from Stata's generator, so \texttt{set seed} before the command makes a fit reproducible. The fits of this manual give the seed in the command so that every listing can be reproduced.

\section{Looking at the fit}

\texttt{estat mcmcdiag} shows how well the algorithm converged:
\showlog{karate_diag}
The column \texttt{t} is the convergence statistic of each coefficient. With stochastic approximation the procedure repeats itself, up to 20 times, until all are below $0.1$; the header of the fit says whether that happened. \texttt{estat gof} compares the model with simulated outcomes (Chapter~\ref{ch:postest}); \texttt{nwalaam simulate} draws outcomes from the fitted model.

\section{Display}\label{sec:display}

The coefficient names are shown in full. At a line size below 100 the confidence intervals are left out so that names of up to 32 characters fit; \opt{ci}, or an explicit \opt{level(\#)}, shows them, and \opt{noci} leaves them out. \opt{ci} and \opt{noci} together are an error. The other display options of Stata's estimation commands (\opt{nopvalues}, \opt{cformat()} and so on) are passed on. The command without arguments replays the last fit, with these options:
\showlog{glasgow_replay}
The replay needs the preceding fit to be that of \nwalaam{}; \texttt{estimates store} and \texttt{estimates restore} work as for other estimation commands.
